\documentclass[12pt,letterpaper,oneside]{article}
\newcommand{\planCaso}{NexoTeX}
\newcommand{\planDescripcion}{Propuesta hipotética de maquetación y capacitación en LaTeX.}
\newcommand{\planFecha}{2026-09-20}
\newcommand{\planLugar}{Cajeme, Sonora; contexto propuesto}
\input{ITESCA/maestria-en-gestion-administrativa/plan-de-negocios-mga/referencias-plan-de-negocios/notas-plan-de-negocios/materiales-generales/historico-formatos/formato-itesca-plan-negocios-generadas.tex}
\begin{document}
\portadaITESCA{Imagen de la empresa}{6545}
\section{Etapa 1. Diagnóstico de la actividad}

Este documento es un borrador hipotético pendiente de confirmación y adaptación con información real del proyecto. Corresponde a la actividad 6545, denominada \textbf{Imagen de la empresa}, de la asignatura Plan de Negocios, clave curricular GGPN01, de la Maestría en Gestión Administrativa del ITESCA.

La consigna solicita presentar el nombre del proyecto, el logotipo y el lema, así como justificar su relación con la propuesta empresarial. También requiere explicar la facilidad de recuerdo del nombre, la autoría y representación del logotipo, la elección de los colores y la utilidad estratégica del lema. El formato de entrega indicado es Word; sin embargo, no se publicaron la extensión exacta, el tamaño máximo ni el nombre requerido para el archivo.

Para atender la actividad se utiliza el caso provisional \textbf{NexoTeX}, entendido como una propuesta de microservicio de maquetación en LaTeX y capacitación dirigida a docentes, investigadores, estudiantes y pequeñas organizaciones de Cajeme, Sonora. El servicio trabajaría únicamente con documentos redactados por sus propios autores y no ofrecería suplantación académica. No existe todavía una empresa confirmada, validación de mercado, cartera de clientes, infraestructura comercial ni registro de marca.

\section{Etapa 2. Diagnóstico de elaboración, verificación y riesgos}

La elaboración requiere comprobar la congruencia entre el nombre, la identidad visual, el lema y el servicio propuesto. Después deberá integrarse el documento, aplicar la plantilla institucional y convertir el resultado al formato Word solicitado. La revisión final deberá considerar legibilidad, ortografía, correspondencia con la actividad y condiciones de entrega en el aula virtual.

Los criterios básicos de verificación son los siguientes:

\begin{itemize}
\item \textbf{Congruencia institucional:} mantener un tono académico, distinguir las propuestas de los hechos comprobados y respetar el formato institucional que se aplique posteriormente.
\item \textbf{Pertinencia curricular:} relacionar la identidad con una propuesta de negocio adecuada para la asignatura Plan de Negocios.
\item \textbf{Correspondencia con la actividad:} incluir nombre, logotipo y lema, además de justificar significado, representación, colores, autoría y utilidad estratégica.
\item \textbf{Trazabilidad:} señalar expresamente los elementos que todavía requieren validación, en lugar de presentarlos como resultados definitivos.
\end{itemize}

Los principales riesgos identificados son la posible coincidencia del nombre con una marca existente, la falta de pruebas de recordación con usuarios y la reproducción inconsistente de los colores. También está pendiente confirmar la apertura de la actividad, pues la referencia ``2006'' se conserva tal como aparece y requiere aclaración. El valor 30 asignado a la unidad 2 no debe interpretarse como la ponderación individual de esta tarea. Otros riesgos son la ausencia del archivo de imagen, la existencia de una rúbrica todavía no consultada, el incumplimiento del formato Word o el uso de un nombre de archivo diferente al solicitado.

\section{Etapa 3. Propuesta de imagen de la empresa}

\subsection{Nombre del proyecto}

El nombre propuesto es \textbf{NexoTeX}. La palabra ``nexo'' comunica la idea de vínculo entre el contenido redactado por una persona y su presentación documental. El componente ``TeX'' alude al sistema tipográfico en el que se apoyaría el servicio para estructurar documentos técnicos, académicos y organizacionales. En conjunto, el nombre busca expresar la conexión entre contenido, organización y edición.

La elección resulta congruente con un microservicio que ofrecería apoyo de maquetación y capacitación. No se propone sustituir la autoría intelectual de los usuarios, sino facilitar que sus propios textos sean organizados mediante plantillas editables y criterios consistentes de presentación. Por ello, el término ``nexo'' identifica la función de acompañamiento, mientras que ``TeX'' delimita el recurso tecnológico principal.

Aunque ``TeX'' proviene del nombre de una tecnología y puede no ser familiar para todo el mercado, su inclusión representa una decisión estratégica: diferencia la propuesta y permite asociarla con la especialidad del servicio. La combinación es breve, se pronuncia de manera directa en español y puede conservarse visualmente como una sola unidad. Estas cualidades sugieren una facilidad potencial de recuerdo, pero dicha condición deberá comprobarse con usuarios y no se presenta como resultado de una investigación de mercado.

\subsection{Logotipo propuesto}

\includegraphics[width=0.9\linewidth]{ITESCA/maestria-en-gestion-administrativa/plan-de-negocios-mga/actividades-generadas/2026-II/revision-2026-09-20/marca-nexotex.png}

El logotipo propone una letra N construida mediante dos barras y una diagonal en verde petróleo, acompañada por dos nodos circulares en grafito. El nombre NexoTeX aparece también en grafito. La N relaciona visualmente la inicial del nombre con la idea de conexión. Las barras representan elementos organizados dentro de una estructura documental, mientras que la diagonal establece continuidad entre ellos. Los nodos refuerzan la noción de enlace y edición progresiva.

El verde petróleo corresponde al código propuesto \#126A60. Se seleccionó para distinguir el símbolo principal y proyectar una apariencia sobria que pueda utilizarse en contextos académicos y profesionales. El grafito, con código \#20242A, se asignó a los nodos y al nombre para favorecer el contraste y mantener una lectura neutral. Estos significados son decisiones de diseño y no constituyen afirmaciones psicológicas ni conclusiones derivadas de pruebas con el mercado.

La propuesta gráfica fue generada con asistencia de inteligencia artificial a partir de las características definidas para el proyecto. Debe considerarse un punto de partida sujeto a revisión humana. Antes de adoptarla se requiere comprobar su originalidad visual, disponibilidad marcaria, legibilidad en tamaños pequeños, reproducción en escala de grises y aceptación entre usuarios potenciales. También deberá confirmarse quién asumirá formalmente la autoría, adaptación y responsabilidad sobre la versión definitiva.

\subsection{Lema y función estratégica}

El lema propuesto es \textbf{``Claridad que se puede editar''}. Su origen se encuentra en dos atributos deseados del servicio. ``Claridad'' se refiere a la organización visual y estructural del documento; la expresión ``se puede editar'' destaca que el resultado buscaría mantenerse reutilizable y modificable, en lugar de limitarse a una salida gráfica cerrada.

Estratégicamente, el lema resume la propuesta sin prometer resultados automáticos. Puede utilizarse junto al logotipo en materiales informativos, presentaciones y comunicaciones comerciales. Además, ayuda a explicar que el servicio se orientaría tanto a la presentación clara como a la continuidad de edición por parte del autor. Su comprensión, pertinencia y capacidad de diferenciación todavía deben validarse con personas pertenecientes a los públicos previstos.

\section{Datos necesarios para convertir el borrador en trabajo real}

Antes de la entrega será necesario confirmar la rúbrica, la fecha efectiva de apertura, la extensión, el tamaño y el nombre del archivo. También se deberá verificar la existencia y calidad de la imagen marca-nexotex.png, revisar la disponibilidad del nombre y del signo, identificar al responsable final del diseño y recopilar retroalimentación auténtica de usuarios potenciales.

Finalmente, deberá aplicarse la plantilla institucional, revisar el resultado visual, convertir el documento al formato Word y comprobar manualmente que el archivo abra correctamente. Este borrador no implica que la identidad haya sido validada, que la empresa exista formalmente ni que la actividad haya sido enviada al aula virtual.
\clearpage
\printbibliography[title={Fuentes de consulta}]
\end{document}
